\chapter{Extreme values, hitting times and a path to RH}\label{chap:evt}\label{ch:evt}\chaptermark{Extreme values and hitting times}

The P\'olya density whose cosine transform is the Riemann $\Xi$ function is an exact superposition of extreme value laws, and the reciprocal $\xi$ function is the Laplace transform of a sum of independent exponentials whose rates are the squared ordinates $\gamma_k^2$. Kent's theorem reads such a sum as a hitting time of a one dimensional diffusion, and Krein's inverse spectral theorem says that every strictly increasing sequence of positive rates arises this way. This chapter puts the two halves together. Extreme value theory supplies explicit approximations to the P\'olya density, each approximation defines a candidate hitting time, and the question of real and simple zeros becomes the question of whether that hitting time belongs to a string with positive mass. The chapter asks whether this gives a path to RH. It does not prove RH. It gives a precise target, three routes to it with the place where each one stops, a family of approximations (the folded Gumbel kernels) whose zeros can be tracked rigorously to a stated height, and a set of diagnostics, from Hankel determinants of the kernel moments to threshold fits of the hitting time tail, whose power and limits are measured.

The chapter develops a source chapter on extreme values, hitting times and a path to RH by Nick Polson \cite{P26evt}, together with the computations done for this book. Its numbers were recomputed here by scripts written from the mathematics alone (the scripts of the source were not used; the scripts of this chapter are \texttt{code/ev\_g1\_*.py}, \texttt{code/ev\_evt\_lib.py} and the other files \texttt{code/ev\_*.py}, listed in \texttt{code/README.md}). Every number printed below is one of those recomputations unless it is marked unchecked, and where the recomputation differed from the source the recomputed value is printed. Statements of papers that were not read are marked unchecked in the text. The notation is the book's, with the following warnings. The kernel $\Phi$ of this chapter is the kernel $\Phi_{\mathrm{dBN}}$ of Section~\ref{ch:dbn} (the normalisation of Rodgers and Tao, with $e^{9u}$ and $e^{4u}$), not the kernel $\Phi$ of \eqref{eq:ch2:Phi}; with it $H(z)=\int_0^\infty\Phi(u)\cos(zu)\,du=\tfrac18\Xi(z/2)$, so the zeros of $H$ are at $z=\pm2\gamma_k$ and the variable of the zeros is $x=z/2$. The eigenvalues of a string are $\lambda_k$, the time of the de Bruijn and Newman flow is $t$ as in Section~\ref{ch:dbn} and $\Lambda$ is its constant, and Euler's constant is $\gamma_E$.

Several parts of the source are already in the book and are cited, not restated: Kent's theorem and the reduction of RH to a string with simple spectrum (Proposition~\ref{ht:prop:kent}, Proposition~\ref{ht:prop:red}, Section~\ref{sec:ch18:kent}), the B spline form of the truncated hitting time and its relation to the Gumbel law (Section~\ref{ht:sec:ev}, Proposition~\ref{ht:prop:bspline}), the hazard rate and the Prony extraction of the atoms (Sections~\ref{sec:ch18:law} and~\ref{sec:ch18:reach}), the Liouville operator and the explicit diffusion for P\'olya's kernel (Proposition~\ref{prop:biane:liouville} and Theorem~\ref{thm:biane:clock}), the Stieltjes criterion for the logarithmic derivative (Theorem~\ref{mw:thm:stieltjes}), the incomplete gamma form of the truncated theta series, the kink at the origin and the tracking of zeros by Rouch\'e's theorem (Proposition~\ref{mi:prop:incgamma}, Theorem~\ref{mi:thm:kink}, Corollary~\ref{mi:cor:rouche}) and the de Bruijn and Newman flow (Section~\ref{ch:dbn}). What is new here is the decomposition of the P\'olya density into Gumbel pieces, the sharp cutoff (folded) family and its height rule, the Hankel string test, the shifted Bessel strings, the detection of the first nonreal pair and the diagnostics of the hitting time tail.

Status words follow the rest of the book. A statement is proved when a proof or a pointer is given, computed when a script of the book produced it, certified when it was computed in ball arithmetic and the sentence says what the certificate covers, conditional, heuristic or open as stated. A computation in ordinary floating point or on a sampled set of points is called computed and never certified. The results are these.
\begin{enumerate}
\item Each term of the theta series for $\Phi$ is an explicit combination of two Gumbel kernels of fractional order, and the Mellin transform of the term is $n^{-s}\xi(s)/(4\zeta(s))$ (Theorem~\ref{ev:thm:decomp}, proved). This gives three families of approximations: the folded Gumbel kernels, P\'olya's two sided kernel and the Dyck chain. Their real zeros, to the heights tabulated, are computed (Table~\ref{ev:tab:approx}). The height $T_N\approx4(N+1)^2$ at which the folded family fails is a heuristic balance of the kink against the size of $H$.
\item The zeros of $\Xi$ are all real and simple if and only if the reciprocal $\xi$ law is the hitting time of a string (Corollary~\ref{ev:cor:rh}, proved, with the hypotheses of simplicity and of the class of $\widehat H$ that the source left out). The hitting time tail gives the principal eigenvalue $\gamma_1^2$ without any zero, and a Laplace inversion of the tail recovers $\gamma_1$ to eighteen digits (Proposition~\ref{ev:prop:tail}, proved under RH, simple zeros and the lower bound $|\xi'(\rho)|\ge e^{-C|\gamma|}$ for the full series, the exponential excess needing only the two term tail; the digits are computed).
\item The Gumbel pieces have unconditional Hadamard products whose atoms are the poles of the archimedean factor, while the atoms $\gamma_k^2$ come from the superposition, which is a sum and not a product (Section~\ref{ev:sec:bspline}, proved for the pieces and heuristic for the superposition).
\item P\'olya's kernel has an explicit string, the Liouville string of Biane, and its first eigenvalues are computed (Section~\ref{ev:sec:liouville}).
\item For any candidate kernel with infinitely many zeros the zeros are all real if and only if the Hankel determinants of the moment sequence $\mu_j=\sum_k\widehat\gamma_k^{-2j}$ are all positive (Proposition~\ref{ev:prop:hankel}, proved). The first eleven pairs of determinants are positive for the exact kernel, the two term folded kernel and P\'olya's kernel (computed). Finite order positivity cannot see an off line pair above a height that grows with the order (heuristic).
\item There are three routes to RH from this material; each stops at a known place, and the folded family gives a reduction of RH to a height rule for all $N$ (Section~\ref{ev:sec:path}, Proposition~\ref{ev:prop:reduction}, proved). The Rouch\'e inequality behind the height rule is certified in ball arithmetic for $N=1,2,3,4$ and the heights of the table, given the count of zeros of $\Xi$ below the height (Section~\ref{ev:sec:work}); the extension to $N=5$ is computed.
\item The eigenvalue problem for the shifted Bessel kernels is a rational recurrence, and its eigenvalues are real (proved), but the positivity of the extension and its reading as a string are not shown for the tabulated cases; the cotangent law for fractional shifts, the collision of the de Bruijn and Newman zeros and the masking of the first nonreal pair by finite Hankel orders are treated in Sections~\ref{ev:sec:work} and~\ref{ev:sec:round2} with the status stated there.
\item Devroye's alternating series gives a sign test for the reciprocal $\Xi$ law, the mean excess function of the hitting time has an expansion with a rate, and a detection rule for the first nonreal pair in the de Bruijn and Newman family is stated and tested (Section~\ref{ev:sec:devroye}, computed unless stated).
\end{enumerate}
The remaining open problems close the chapter.
